\documentclass[11pt,a4paper]{article}
\usepackage[T1]{fontenc}
\usepackage[utf8]{inputenc}
\usepackage{lmodern,amsmath,amssymb}
\usepackage[margin=25mm]{geometry}
\usepackage[hidelinks]{hyperref}
\hypersetup{pdftitle={The Agmon bound controls the global excess only: the local large-field estimate does not f},pdfauthor={navstokgap research working notes}}
\newcommand{\tightlist}{\setlength{\itemsep}{0pt}\setlength{\parskip}{0pt}}
\setlength{\emergencystretch}{3em}
\setcounter{secnumdepth}{0}
\begin{document}
\hypertarget{the-agmon-bound-controls-the-global-excess-only-the-local-large-field-estimate-does-not-follow}{%
\section{The Agmon bound controls the global excess only: the local
large-field estimate does not
follow}\label{the-agmon-bound-controls-the-global-excess-only-the-local-large-field-estimate-does-not-follow}}

Carrying out the distance estimate of
\href{agmon-ground-state-suppression.md}{the Agmon note} §2 with
constants changes its meaning. The rigorous chain is
\[d(U)\ \ge\ \sqrt{\frac{B'}{A'}}\int_0^{f_0}\frac{\sqrt f\,df}{\|\nabla V\|_\infty}
=\frac{2}{3}\,\frac{2}{g^2}\,\frac{f_0^{3/2}}{\|\nabla V\|_\infty},
\qquad f_0=V(U)-\bar v,\] obtained from
\(|df|\le\|\nabla V\|_\infty|d\gamma|\) along any path to the allowed
region, and the gradient norm of the magnetic energy over the whole
configuration space is
\[\|\nabla V\|_\infty\ \le\ 2\sqrt{2N|\mathcal E|},\] \textbf{extensive
in the volume}. For an excess \(w\) per plaquette spread over the
lattice this gives an extensive exponent,
\[\big|\Omega\big|^2\ \lesssim\ \exp\Big[-\frac{2\sqrt2}{3}\,\frac{w^{3/2}}{g^2\sqrt N}\,|\mathcal P|\Big],\]
with the same rate per plaquette as the earlier estimate when
\(w=O(N)\). For an excess carried by \(n\) plaquettes in a \textbf{fixed
region} of a large lattice the same chain gives
\(d\gtrsim n^{3/2}N/(g^2\sqrt{|\mathcal P|})\), which tends to zero as
the volume grows: the argument controls the total magnetic energy and
says nothing about a local region. The reason is structural rather than
technical. The forbidden region of the Agmon method is defined by the
\textbf{total} potential exceeding the \textbf{total} energy, and in an
extensive system a local excess never makes the total exceed its
extensive mean. So \href{agmon-ground-state-suppression.md}{the Agmon
note} proves a thermodynamic large-deviation bound, and the local
large-field estimate that the renormalization step needs does not follow
from it. The gap between the two is the same one the Euclidean
formulation closes for free, because there the weight \(e^{-S_w}\)
factorizes over plaquettes. Constants explicit; this note corrects the
previous one; nothing promoted.

\hypertarget{the-rigorous-distance-bound}{%
\subsection{1. The rigorous distance
bound}\label{the-rigorous-distance-bound}}

Let \(\bar v=e_0/B'\) and \(f=V-\bar v\), so the allowed region is
\(\{f\le0\}\) and the Agmon metric is \(\sqrt{(B'/A')f}\,|dU|\) on
\(\{f>0\}\).

\textbf{Proposition 1.} For every \(U\) with \(f_0=f(U)>0\),
\[d(U)\ \ge\ \frac{2}{3}\sqrt{\frac{B'}{A'}}\;\frac{f_0^{3/2}}{\|\nabla V\|_\infty}
=\frac{4}{3g^2}\;\frac{f_0^{3/2}}{\|\nabla V\|_\infty}.\]

\emph{Proof.} Let \(\gamma\) be any path from \(U\) to \(\{f\le0\}\),
parametrized by arclength. Along it
\(|f'(s)|\le|\nabla V(\gamma(s))|\le\|\nabla V\|_\infty\), so
\[\int_\gamma\sqrt{\tfrac{B'}{A'}f}\;ds
\ \ge\ \sqrt{\tfrac{B'}{A'}}\int_\gamma\sqrt{f}\;\frac{|f'|\,ds}{\|\nabla V\|_\infty}
\ \ge\ \frac{\sqrt{B'/A'}}{\|\nabla V\|_\infty}\int_0^{f_0}\sqrt f\,df
=\frac{2}{3}\frac{\sqrt{B'/A'}}{\|\nabla V\|_\infty}f_0^{3/2},\] using
that \(f\) runs from \(f_0\) to \(0\). Take the infimum over \(\gamma\)
and insert \(\sqrt{B'/A'}=2/g^2\). \(\square\)

\textbf{Lemma 2.} \(\|\nabla V\|_\infty\le2\sqrt{2N|\mathcal E|}\) for
\(SU(N)\) in the normalization of
\href{mass-gap-obligations-lattice.md}{the obligations map}.

\emph{Proof.} For one link,
\(|\nabla_\ell V|\le\sum_{p\ni\ell}|\nabla_\ell\operatorname{Re}\operatorname{tr}U_p|\),
and the completeness relation of \href{lattice-gap-upper-bounds.md}{the
upper-bound note} Corollary 2 gives
\(|\nabla_\ell\operatorname{Re}\operatorname{tr}U_p|\le\sqrt{N/2}\);
each link lies in four plaquettes, so
\(|\nabla_\ell V|\le4\sqrt{N/2}=2\sqrt{2N}\). Summing the squares over
the \(|\mathcal E|\) links gives \(\|\nabla V\|^2\le8N|\mathcal E|\).
\(\square\)

\hypertarget{what-follows-and-what-does-not}{%
\subsection{2. What follows, and what does
not}\label{what-follows-and-what-does-not}}

\emph{Global excess.} With \(f_0=w|\mathcal P|\) and
\(|\mathcal E|=|\mathcal P|\) in three dimensions,
\[d\ \ge\ \frac{4}{3g^2}\,\frac{(w|\mathcal P|)^{3/2}}{2\sqrt{2N|\mathcal P|}}
=\frac{\sqrt2}{3}\,\frac{w^{3/2}}{g^2\sqrt N}\,|\mathcal P| ,\]
extensive, and the Agmon estimate of
\href{agmon-ground-state-suppression.md}{the Agmon note} Corollary 2
gives \[\int_{\{V-\bar v\ge w|\mathcal P|\}}\Omega^2\,d\mu
\ \le\ C\exp\Big[-\frac{2\sqrt2(1-\delta)}{3}\,\frac{w^{3/2}}{g^2\sqrt N}\,|\mathcal P|\Big].\]
At \(w=O(N)\) the rate per plaquette is of order \(N/g^2\), as the
previous note stated.

\emph{Local excess.} If instead the excess \(nv\) sits on \(n\)
plaquettes of a fixed region while the rest of the lattice is at its
mean, then \(f_0=nv\) and
\[d\ \ge\ \frac{4}{3g^2}\,\frac{(nv)^{3/2}}{2\sqrt{2N|\mathcal P|}}
\ \sim\ \frac{n^{3/2}N}{g^2\sqrt{|\mathcal P|}}\ \xrightarrow[\ |\mathcal P|\to\infty\ ]{}\ 0 .\]
The bound degrades with the volume and gives nothing in the
thermodynamic limit.

\hypertarget{why-the-local-statement-fails-structurally}{%
\subsection{3. Why the local statement fails
structurally}\label{why-the-local-statement-fails-structurally}}

The Agmon method measures tunnelling into the region where the
\textbf{total} potential exceeds the \textbf{total} energy. In an
extensive system \(\bar v=e_0/B'\) is extensive, of order
\(N|\mathcal P|\), so a local excess of \(n\) plaquettes leaves \(V\)
far below \(\bar v\) whenever the rest of the lattice sits at or below
its mean. Such a configuration is classically allowed, the Agmon weight
vanishes on it, and no decay is asserted.

A localized version would need a weight \(\rho\) supported near the
region, subject to \(A'|\nabla\rho|^2\le B'V-e_0\) \textbf{pointwise},
and that inequality fails for exactly the same reason: its right side is
a global quantity that a local excess does not control. Repairing it
requires a local energy balance, that is, a statement that the
ground-state energy density is locally attained, which is a form of the
very locality one is trying to prove.

\hypertarget{consequence-for-the-route}{%
\subsection{4. Consequence for the
route}\label{consequence-for-the-route}}

\emph{What survives.} A thermodynamic large-deviation bound on the total
magnetic energy in the ground state, with the coupling dependence
\(1/g^2\) and an extensive rate. It is a genuine property of the
Kogut--Susskind ground state and it is proved with explicit constants
above.

\emph{What was claimed too strongly.}
\href{agmon-ground-state-suppression.md}{The Agmon note} read the same
estimate as a suppression of a \textbf{local} large-field region, which
is what the renormalization step needs, and that reading does not
follow. Its Sections 3 and 4 should be read with ``global excess'' in
place of ``\(n\) excited plaquettes'', and its comparison with the
Euclidean Wilson weight holds only for the global statement. The
Euclidean weight \(e^{-S_w}\) factorizes over plaquettes and therefore
gives the local statement immediately, which is the structural advantage
the Hamiltonian formulation lacks and which
\href{typical-field-strength-window.md}{the typical-field note} had
already identified.

\emph{What would close it.} Either a local energy-balance estimate for
the Kogut--Susskind ground state, or a proof that the ground-state
measure is dominated by a product-form weight,
\(\Omega^2\le Ce^{-\lambda V}\) with \(\lambda>0\), which is a statement
of exactly the kind the strong-coupling expansion produces and which is
open at intermediate coupling. The second is the sharper target: it is a
single inequality, it implies the local estimate by factorization, and
at \(g=\infty\) it holds with \(\lambda=0\).

\hypertarget{consequence-for-state}{%
\subsection{5. Consequence for STATE}\label{consequence-for-state}}

The Agmon route yields the global large-deviation bound and stops there.
The local estimate, which is what the decimation step requires, needs a
pointwise domination of the ground state by a Gibbs weight,
\(\Omega^2\le Ce^{-\lambda V}\) with \(\lambda>0\) uniform in the
volume. That is now the single sharpest open question in this line, it
is elementary to state, and the strong-coupling case should be provable
by perturbation theory around \(\Omega\equiv1\).
\end{document}
